\documentclass[9pt,landscape,a4paper]{extarticle}
\usepackage[utf8]{inputenc}
\usepackage[T1]{fontenc}
\usepackage[sfdefault]{noto-sans}  % Noto Sans font
\usepackage[ttdefault=true]{sourcecodepro}  % Source Code Pro for monospace (straight quotes)
\usepackage{upquote}  % Straight quotes in verbatim/ttfamily
\usepackage{tikz}
\usetikzlibrary{positioning,arrows.meta,shapes.geometric,fit}
\usepackage[landscape,margin=0.3cm]{geometry}
\usepackage{multicol}
\usepackage{amsmath}
\usepackage{array}
\usepackage{enumitem}
\usepackage{xcolor}
\usepackage{titlesec}
\usepackage{fancyhdr}
\usepackage[hidelinks]{hyperref}
\usepackage{tcolorbox}
\tcbuselibrary{listings,skins}

% Formatting with balanced spacing
\setlength{\parindent}{0pt}
\setlength{\parskip}{0pt}
\setlength{\columnsep}{0.4cm}
\setlist[itemize]{noitemsep, topsep=0pt, leftmargin=1.2em, parsep=0pt}

% Reduce verbatim spacing
\usepackage{etoolbox}
\makeatletter
\preto{\@verbatim}{\topsep=0pt \partopsep=0pt}
\makeatother

% Colors
\definecolor{sectioncolor}{RGB}{0,102,204}
\definecolor{codebg}{RGB}{245,245,245}
\definecolor{codeframe}{RGB}{200,200,200}
\definecolor{tipbg}{RGB}{232,245,233}
\definecolor{tipframe}{RGB}{76,175,80}

% Code box style
\newtcolorbox{codebox}{
  colback=codebg,
  colframe=codeframe,
  boxrule=0.3pt,
  arc=1pt,
  left=2pt,right=2pt,top=1pt,bottom=1pt,
  fontupper=\ttfamily\scriptsize,
  before skip=1pt,
  after skip=1pt
}

% Tip/note box style
\newtcolorbox{tipbox}{
  colback=tipbg,
  colframe=tipframe,
  boxrule=0.3pt,
  arc=1pt,
  left=2pt,right=2pt,top=1pt,bottom=1pt,
  fontupper=\scriptsize,
  before skip=1pt,
  after skip=1pt
}

% Section formatting
\titleformat{\section}{\normalfont\bfseries\small\color{sectioncolor}}{\thesection}{0.3em}{}
\titlespacing*{\section}{0pt}{0.5ex}{0.1ex}

% Header colors
\definecolor{headerblue}{RGB}{0,102,204}

% Header/footer
\pagestyle{fancy}
\fancyhf{}
\renewcommand{\headrulewidth}{0pt}
\fancyfoot[C]{\small LINE v3.0 -- \url{http://line-solver.sf.net}}

\begin{document}

% Title
\begin{center}
{\Large\bfseries\color{headerblue} LINE Solver Cheatsheet}\\[-0.2em]
{\small\url{http://line-solver.sf.net}}
\end{center}

\vspace{-0.2em}

% Model Structure - full width box with 4 language columns aligned to cheatsheet columns
\newlength{\colwidth}
\setlength{\colwidth}{0.2425\textwidth}
\begin{tcolorbox}[colback=codebg, colframe=codeframe, boxrule=0.3pt, arc=1pt,
  left=2pt, right=2pt, top=1pt, bottom=1pt, title={\small\textbf{Model Structure}},
  coltitle=black, colbacktitle=codebg]
\ttfamily\tiny\setlength{\tabcolsep}{0pt}
\begin{tabular*}{\textwidth}{@{\extracolsep{\fill}}p{\colwidth}p{\colwidth}@{\hspace{5mm}}p{\colwidth}@{\hspace{5mm}}p{\colwidth}@{}}
\textbf{\textsf{\scriptsize MATLAB}} & \textbf{\textsf{\scriptsize Java}} & \textbf{\textsf{\scriptsize Kotlin}} & \textbf{\textsf{\scriptsize Python}} \\[1pt]
model = Network('name'); &
model = new Network("name"); &
val model = Network("name") &
model = Network("name") \\
source = Source(model,'Src'); &
source = new Source(model,"Src"); &
val source = Source(model,"Src") &
source = Source(model,"Src") \\
queue = Queue(model,'Q', &
queue = new Queue(model,"Q", &
val queue = Queue(model,"Q", &
queue = Queue(model,"Q", \\
\ \ SchedStrategy.FCFS); &
\ \ SchedStrategy.FCFS); &
\ \ SchedStrategy.FCFS) &
\ \ SchedStrategy.FCFS) \\
sink = Sink(model,'Snk'); &
sink = new Sink(model,"Snk"); &
val sink = Sink(model,"Snk") &
sink = Sink(model,"Snk") \\
oclass = OpenClass(model,'C1'); &
oclass = new OpenClass(model,"C1"); &
val oclass = OpenClass(model,"C1") &
oclass = OpenClass(model,"C1") \\
source.setArrival(oclass,Exp(1)); &
source.setArrival(oclass,new Exp(1)); &
source.setArrival(oclass,Exp(1.0)) &
source.setArrival(oclass,Exp(1)) \\
queue.setService(oclass,Exp(2)); &
queue.setService(oclass,new Exp(2)); &
queue.setService(oclass,Exp(2.0)) &
queue.setService(oclass,Exp(2)) \\
model.link(Network.serialRouting( &
model.link(Network.serialRouting( &
model.link(Network.serialRouting( &
model.link(Network.serialRouting( \\
\ \ source,queue,sink)); &
\ \ source,queue,sink)); &
\ \ source,queue,sink)) &
\ \ source,queue,sink)) \\
MVA(model).avgTable() &
new SolverMVA(model).getAvgTable() &
SolverMVA(model).avgTable() &
SolverMVA(model).getAvgTable() \\
\end{tabular*}
\end{tcolorbox}

\vspace{0.1em}

\begin{multicols}{4}

\section*{Node Types}
\begin{tabular}{@{}ll@{}}
\texttt{Cache} & Caching station \\
\texttt{ClassSwitch} & Class switching \\
\texttt{Delay} & Infinite server \\
\texttt{Fork} & Forks jobs \\
\texttt{Join} & Joins jobs \\
\texttt{Place} & Petri net place \\
\texttt{Queue} & Queueing station \\
\texttt{Router} & Routing node \\
\texttt{Sink} & Job departures (open) \\
\texttt{Source} & Job arrivals (open) \\
\texttt{Transition} & Petri net transition \\
\end{tabular}

\section*{Job Classes}
\texttt{ClosedClass(model, 'name', N, refStat)}\\
\texttt{OpenClass(model, 'name')}\\
\texttt{SelfLoopingClass(model, 'name', N, ref)}

\begin{tipbox}
\textbf{Tip:} The reference station (refStat/ref) is where response time and throughput are measured for closed classes. For SelfLoopingClass, jobs remain at the reference station.
\end{tipbox}

\section*{Distributions}
\begin{tabular}{@{}ll@{}}
\texttt{APH($\boldsymbol{\alpha}$, T)} & Acyclic PH \\
\texttt{Coxian($\boldsymbol{\mu}$, $\boldsymbol{\phi}$)} & Coxian \\
\texttt{Det(v)} & Deterministic \\
\texttt{Disabled} & None \\
\texttt{Erlang($\lambda$, k)} & Erlang-k \\
\texttt{Exp($\lambda$)} & Exponential \\
\texttt{Gamma($\alpha$, $\beta$)} & Gamma \\
\texttt{HyperExp(p, $\lambda_1$, $\lambda_2$)} & Hyper-exp \\
\texttt{Lognormal($\mu$, $\sigma$)} & Log-normal \\
\texttt{MAP(D0, D1)} & MAP \\
\texttt{MMPP2($\lambda_1$, $\lambda_2$, $\sigma_1$, $\sigma_2$)} & MMPP(2) \\
\texttt{Pareto($\alpha$, k)} & Pareto \\
\texttt{PH($\boldsymbol{\alpha}$, T)} & Phase-type \\
\texttt{Replayer(trace)} & Trace replay \\
\texttt{Uniform(a, b)} & Uniform \\
\texttt{Weibull($\alpha$, $\beta$)} & Weibull \\
\end{tabular}

\begin{tipbox}
\textbf{Tip:} Fit distributions to moments: \texttt{Exp.fitMean(m)}, \texttt{Erlang.fitMeanAndSCV(m,scv)}, \texttt{APH.fit(mean,scv,skew)}.
\end{tipbox}

\section*{Scheduling Strategies}
\begin{tabular}{@{}ll@{}}
\texttt{DPS} & Discriminatory PS \\
\texttt{DPSPRIO} & DPS with Priority \\
\texttt{FB} & Feedback (LAS) \\
\texttt{FCFS} & First-Come First-Served \\
\texttt{FCFSPRIO} & FCFS with Priority \\
\texttt{GPS} & Generalized PS \\
\texttt{GPSPRIO} & GPS with Priority \\
\texttt{HOL} & Head-of-Line Priority \\
\texttt{INF} & Infinite Servers \\
\texttt{LCFS} & Last-Come First-Served \\
\texttt{LCFSPRIO} & LCFS with Priority \\
\texttt{LCFSPR} & LCFS Preemptive \\
\texttt{LEPT} & Longest Expected PT \\
\texttt{LJF} & Longest Job First \\
\texttt{LRPT} & Longest Remaining PT \\
\texttt{PS} & Processor Sharing \\
\texttt{PSJF} & Preemptive SJF \\
\texttt{PSPRIO} & PS with Priority \\
\texttt{SEPT} & Shortest Expected PT \\
\texttt{SETF} & Shortest Elapsed Time First \\
\texttt{SIRO} & Random (SIRO) \\
\texttt{SJF} & Shortest Job First \\
\texttt{SRPT} & Shortest Remaining PT \\
\end{tabular}

\begin{tipbox}
\textbf{Tip:} Set DPS/GPS weights with \texttt{queue.setService(class, dist, weight)}. Set priorities via \texttt{ClosedClass(..., priority)} or \texttt{class.setPriority(p)}. Only \texttt{*PRIO} policies are sensitive to priorities.
\end{tipbox}

\section*{Solvers}
\begin{tabular}{@{}ll@{}}
\texttt{AUTO} & Auto-select best \\
\texttt{CTMC} & CTMC (exact, small) \\
\texttt{LDES} & Discrete Event Sim \\
\texttt{ENV} & Random environments \\
\texttt{FLD} & Fluid/mean-field ODEs \\
\texttt{JMT} & JMT simulator \\
\texttt{LN} & Layered networks \\
\texttt{LQNS} & LQNS interface \\
\texttt{MAM} & Matrix Analytic Methods \\
\texttt{MVA} & Mean Value Analysis \\
\texttt{NC} & Normalizing Constant Methods \\
\texttt{QNS} & LQNS product-form solver \\
\texttt{SSA} & Stochastic simulation \\
\end{tabular}

\begin{tipbox}
\textbf{Tip:} Use \texttt{listMethods()} to see available methods for each solver. Use \texttt{MVA} for product-form networks (fastest), \texttt{CTMC} for exact small models, \texttt{JMT}/\texttt{SSA} for general simulation.
\end{tipbox}

\section*{Solver Options}
\begin{codebox}
MVA(model, 'method', 'exact')\\
JMT(model, 'seed', 123,\\
\quad 'samples', 10000)\\
SSA(model, 'samples', 5000,\\
\quad 'method', 'para')
\end{codebox}

\section*{Output Metrics}
\begin{tabular}{@{}ll@{}}
\texttt{ArvR} & Arrival rate \\
\texttt{QLen} & Mean queue length \\
\texttt{ResidT} & Residence time (total across visits) \\
\texttt{RespT} & Response time (per visit) \\
\texttt{Tput} & Throughput \\
\texttt{Util} & Server utilization \\
\end{tabular}

\section*{Analysis Methods}
\begin{codebox}
solver.avgTable()  \% alias: aT\\
solver.avgQLen(), avgUtil()\\
solver.avgRespT(), avgTput()\\
solver.avgSysRespT()\\
solver.avgSysTput()\\
solver.avgChainQLen()
\end{codebox}

\section*{Method Aliases}
\begin{tabular}{@{}ll@{}}
\texttt{aCT} & \texttt{getAvgChainTable} \\
\texttt{aNT} & \texttt{getAvgNodeTable} \\
\texttt{aST} & \texttt{getAvgSysTable} \\
\texttt{aT} & \texttt{getAvgTable} \\
\end{tabular}

\section*{Routing}
\begin{codebox}
model.link(Network.serialRouting(\\
\quad n1,n2,n3));\\
n1.setProbRouting(class, n2, 0.7);\\
n1.setRouting(class,\\
\quad RoutingStrategy.RROBIN);
\end{codebox}

\section*{Multi-Server \& Finite Capacity}
\begin{codebox}
queue.setNumServers(k);\\
queue.setCapacity(bufferSize);
\end{codebox}

\section*{Class Switching}
\begin{codebox}
cs = ClassSwitch(model, 'CS');\\
P = [0.3 0.7; 0.5 0.5];\\
cs.setClassSwitching(P);
\end{codebox}

\section*{Fork-Join Networks}
\begin{codebox}
fork = Fork(model, 'Fork');\\
join = Join(model, 'Join', fork);
\end{codebox}

\section*{Layered Networks}
\begin{codebox}
lqn = LayeredNetwork('name');\\
P = Processor(lqn,'P',1,\\
\quad SchedStrategy.PS);\\
T = Task(lqn,'T',1,\\
\quad SchedStrategy.REF).on(P);\\
E = Entry(lqn,'E').on(T);\\
A = Activity(lqn,'A').on(T)\\
\quad .boundTo(E).setHostDemand(Exp(1));
\end{codebox}

\section*{Caching}
\begin{codebox}
cache = Cache(model,'Cache',\\
\quad nItems,cap,ReplacementStrategy.LRU);\\
cache.setPopularity(Zipf(1.4));
\end{codebox}

\section*{Model Import/Export}
\begin{codebox}
model.view()  \% Open in JMT\\
model.save('file.jsimg')\\
Network.load('file.jsimg')
\end{codebox}

\begin{tipbox}
\textbf{Tip:} Check model properties: \texttt{model.hasProductFormSolution()} (exact MVA possible), \texttt{model.getNumStations()}, \texttt{model.getNumClasses()}, \texttt{model.getStruct()} (internal representation).
\end{tipbox}

\section*{Routing Strategies}
\begin{tabular}{@{}ll@{}}
\texttt{JSQ} & Join shortest queue \\
\texttt{KCHOICES} & Power of k choices \\
\texttt{PROB} & Probabilistic routing \\
\texttt{RAND} & Random selection \\
\texttt{RROBIN} & Round-robin \\
\end{tabular}

\end{multicols}

% QN diagram at bottom right - outside multicols (COMMENTED OUT)
% \begin{tikzpicture}[remember picture, overlay]
% \node[anchor=south east, xshift=-0.3cm, yshift=0.5cm] at (current page.south east) {
% \begin{tikzpicture}[
%   xscale=0.84,
%   yscale=0.87,
%   % Disk queues and servers (blue)
%   queue/.style={draw, thick, fill=headerblue!15, minimum width=0.5cm, minimum height=0.25cm, rounded corners=1pt},
%   server/.style={draw, thick, fill=headerblue!30, circle, minimum size=0.2cm},
%   % CPU queue and server (orange/red)
%   cpuqueue/.style={draw, thick, fill=orange!20, minimum width=0.5cm, minimum height=0.25cm, rounded corners=1pt},
%   cpuserver/.style={draw, thick, fill=orange!40, circle, minimum size=0.2cm},
%   % Delay servers (green)
%   delayserver/.style={draw, semithick, fill=green!30, circle, minimum size=0.03cm},
%   arr/.style={-{Stealth[length=1.5mm]}, semithick, headerblue!70},
% ]
%   % Delay/Think station (left) - two vertical lines with 3 stacked circles (green), shifted up 5mm
%   \draw[thick, green!50!black] (-3.25,0.05) -- (-3.25,0.95);
%   \draw[thick, green!50!black] (-2.55,0.05) -- (-2.55,0.95);
%   \node[delayserver] (t1) at (-2.80,0.83) {};
%   \node[delayserver] (t2) at (-2.80,0.5) {};
%   \node[delayserver] (t3) at (-2.80,0.17) {};
%   \coordinate (think-in) at (-2.80,0.05);
%   \coordinate (think-out) at (-2.55,0.5);
%   \node[above=0.12cm of t1, font=\tiny, green!50!black] {Delay};
%
%   % Central CPU (orange)
%   \node[cpuqueue] (cpuq) at (0,0) {};
%   \node[cpuserver, right=0.05cm of cpuq] (cpus) {};
%   \node[below=0.15cm of cpuq, font=\tiny\bfseries, orange!70!black] {CPU};
%
%   % Disk 1 (top right) - blue
%   \node[queue] (d1q) at (3.2,1.3) {};
%   \node[server, right=0.05cm of d1q] (d1s) {};
%   \node[below=0.15cm of d1q, font=\tiny, headerblue] {D1};
%
%   % Disk 2 (right) - blue - aligned with D1, D3
%   \node[queue] (d2q) at (3.2,0) {};
%   \node[server, right=0.05cm of d2q] (d2s) {};
%   \node[below=0.15cm of d2q, font=\tiny, headerblue] {D2};
%
%   % Disk 3 (bottom right) - blue
%   \node[queue] (d3q) at (3.2,-1.3) {};
%   \node[server, right=0.05cm of d3q] (d3s) {};
%   \node[below=0.15cm of d3q, font=\tiny, headerblue] {D3};
%
%   % Arrows - Think to CPU (from middle circle height)
%   \draw[arr] (-2.55,0.5) -- (-0.3,0.5) -- (-0.3,0) -- (cpuq.west);
%
%   % Arrows - CPU to Disks (going above/below to avoid overlap)
%   \draw[arr] (cpus.east) -- ++(0.3,0) |- (d1q.west);
%   \draw[arr] (cpus.east) -- (d2q.west);
%   \draw[arr] (cpus.east) -- ++(0.3,0) |- (d3q.west);
%
%   % Arrows - Disks back to CPU (looping back, entering middle of queue)
%   \draw[arr] (d1s.east) -- ++(0.6,0) |- (-0.15,2.0) -- (-0.15,0.2);
%   \draw[arr] (d2s.east) -- ++(0.9,0) |- (-0.15,-2.3) -- (-0.15,-0.2);
%   \draw[arr] (d3s.east) -- ++(0.6,0) |- (-0.15,-2.0) -- (-0.15,-0.2);
%
%   % Arrow - CPU back to Think (enters from below, going up into middle of delay)
%   \draw[arr] (cpuq.south) -- ++(0,-0.8) -- (-2.80,-0.8) -- (-2.80,0.05);
%
%   % LINE label
%   \node[font=\tiny\bfseries, headerblue] at (1.2,-2.8) {LINE};
% \end{tikzpicture}
% };
% \end{tikzpicture}
\end{document}
